\documentclass[a4paper,11pt]{article}

% ── Pakete ──────────────────────────────────────────────────────────────
\usepackage[T1]{fontenc}
\usepackage[utf8]{inputenc}
\usepackage[ngerman]{babel}
\usepackage{geometry}
\geometry{left=25mm, right=25mm, top=35mm, bottom=28mm, headheight=28pt}

\usepackage{amsmath, amssymb}
\usepackage{booktabs}
\usepackage{array}
\usepackage{longtable}
\usepackage{graphicx}
\usepackage{xcolor}
\usepackage{listings}
\usepackage{hyperref}
\usepackage{enumitem}
\usepackage{titlesec}
\usepackage{fancyhdr}
\usepackage{caption}
\usepackage{float}
\usepackage{microtype}
\usepackage{lmodern}
\usepackage{parskip}
\usepackage{cite}
\usepackage{soul}        % Sperrung (Letterspace)
\usepackage{calc}

% ── CD-Farben (achromat) ─────────────────────────────────────────────────
\definecolor{hsm-black}{RGB}{0,0,0}
\definecolor{hsm-darkgray}{RGB}{51,51,51}
\definecolor{hsm-gray}{RGB}{120,120,120}
\definecolor{hsm-lightgray}{RGB}{235,235,235}
\definecolor{code-bg}{RGB}{245,245,245}
\definecolor{code-kw}{RGB}{30,30,30}
\definecolor{code-str}{RGB}{80,80,80}
\definecolor{code-comment}{RGB}{130,130,130}

% ── Logo-Pfad ────────────────────────────────────────────────────────────
\newcommand{\hsmlogopath}{%
  /Users/philippschafer/Library/CloudStorage/%
  OneDrive-HochschuleMainz-UniversityofAppliedSciences/%
  03_Tragwerkslabor/06_Projekte/2026_03_11_MIB/ORGA/CorporateDesign/%
  Logo-Paket_Hochschule_Mainz_NEU/HSM_Logo_Dachmarke/03 PNG/RGB/%
  HSM_Logo_Dachmarke_RGB_pos.png%
}

% ── Code-Listings ────────────────────────────────────────────────────────
\lstdefinestyle{python}{
  language=Python,
  backgroundcolor=\color{code-bg},
  basicstyle=\ttfamily\footnotesize,
  keywordstyle=\color{hsm-black}\bfseries,
  stringstyle=\color{code-str},
  commentstyle=\color{code-comment}\itshape,
  numbers=left,
  numberstyle=\tiny\color{hsm-gray},
  numbersep=8pt,
  frame=single,
  framerule=0.4pt,
  framesep=4pt,
  breaklines=true,
  breakatwhitespace=true,
  tabsize=4,
  showstringspaces=false,
  captionpos=b,
  xleftmargin=12pt,
  rulecolor=\color{hsm-gray},
}
\lstset{style=python}

% ── Hyperlinks ───────────────────────────────────────────────────────────
\hypersetup{
  colorlinks=true,
  linkcolor=hsm-black,
  urlcolor=hsm-darkgray,
  citecolor=hsm-darkgray,
  pdftitle={IFC → Scene Graph → JaxFEM},
  pdfauthor={Philipp Schäfer, Hochschule Mainz},
}

% ── Kopf-/Fußzeile (CD Briefbogen-Stil) ─────────────────────────────────
\pagestyle{fancy}
\fancyhf{}
% Kopfzeile: links Dokumenttitel, rechts Logo
\fancyhead[L]{%
  \begin{minipage}[b]{0.55\textwidth}
    \footnotesize\color{hsm-darkgray}
    IFC $\to$ Scene Graph $\to$ JaxFEM\\
    \textcolor{hsm-gray}{\small Tragwerkslabor · Hochschule Mainz}
  \end{minipage}%
}
\fancyhead[R]{%
  \includegraphics[height=18pt]{\hsmlogopath}%
}
% Fußzeile: links "University of Applied Sciences", rechts Seitenzahl
\fancyfoot[L]{\footnotesize\color{hsm-gray}University of Applied Sciences}
\fancyfoot[R]{\footnotesize\color{hsm-gray}\thepage}
\renewcommand{\headrulewidth}{0.6pt}
\renewcommand{\footrulewidth}{0.4pt}

% ── Überschriften (CD: fett, Großbuchstaben, Trennlinie darunter) ─────────
\titleformat{\section}
  {\large\bfseries\MakeUppercase}
  {}
  {0em}
  {\thesection\quad}
  [\vspace{2pt}\titlerule]
\titlespacing{\section}{0pt}{18pt}{8pt}

\titleformat{\subsection}
  {\normalsize\bfseries}
  {}
  {0em}
  {\thesubsection\quad}
\titlespacing{\subsection}{0pt}{12pt}{4pt}

% ── Titelseite (Notiz/Protokoll-Stil) ───────────────────────────────────
\newcommand{\hsmtitlepage}{%
  \thispagestyle{empty}
  % Kopfzeile der Titelseite
  \noindent
  \begin{minipage}[t]{0.55\textwidth}
    \vspace{0pt}
    \includegraphics[height=22pt]{\hsmlogopath}
  \end{minipage}%
  \hfill
  \begin{minipage}[t]{0.42\textwidth}
    \vspace{4pt}\raggedleft
    \footnotesize\color{hsm-darkgray}
    \textbf{Philipp Schäfer}\\
    Tragwerkslabor\\
    Fachbereich Technik $|$ Engineering\\[4pt]
    Hochschule Mainz\\
    University of Applied Sciences\\
    Holzstraße 36, 55116 Mainz\\[4pt]
    T~+49\,6131\,628-1366\\
    \href{mailto:philipp.schaefer@hs-mainz.de}{philipp.schaefer@hs-mainz.de}
  \end{minipage}

  \vspace{8pt}
  \noindent\rule{\linewidth}{0.6pt}

  \vspace{18pt}
  \noindent{\LARGE\bfseries\MakeUppercase{IFC $\to$ Scene Graph $\to$ JaxFEM}}

  \vspace{4pt}
  \noindent\rule{\linewidth}{0.4pt}

  \vspace{6pt}
  \noindent\textbf{Meilenstein-Dokumentation $|$ Version 1.0}

  \vspace{4pt}
  \noindent\color{hsm-gray}\small
  Automatische FEM-Berechnung aus IFC-Gebäudemodellen\\
  via semantischen Scene Graphs und NetworkX DiGraph

  \vspace{4pt}
  \noindent\normalcolor\small Mainz, April 2026

  \vspace{16pt}
}

% ═══════════════════════════════════════════════════════════════════════
\begin{document}

\hsmtitlepage

% ── Abstract ────────────────────────────────────────────────────────────
\begin{abstract}
\noindent
Das Notebook \texttt{IFC\_SceneGraph\_FEM.ipynb} realisiert erstmals eine vollautomatische Kette von einem IFC-Gebäudemodell~\cite{ifcopenshell} über einen semantischen Scene Graph bis zu einer FEM-Berechnung mit dem JAX-basierten Solver JaxFEM~\cite{jaxfem}. Randbedingungen, Lastflächen und Materialkennwerte werden automatisch aus der Graph-Topologie und den IFC-Properties abgeleitet. Die Ergebnisse (Verschiebungen, Von-Mises-Spannungen) werden in den Scene Graph zurückgeschrieben und als strukturierter LLM-Prompt exportiert. Der erste vollständige Durchlauf mit 149 IFC-Bauteilen, 1\,958 Tet4-Elementen und 1\,917 Freiheitsgraden war erfolgreich.
\end{abstract}

\tableofcontents
\newpage

% ═══════════════════════════════════════════════════════════════════════
\section{Motivation und Einordnung}

Die manuelle Vorbereitung von FEM-Modellen aus BIM-Daten ist in der Praxis ein erheblicher Aufwand: Geometrieimport, Randbedingungsdefinition, Netzerstellung und Materialzuweisung werden heute in RFEM, ANSYS oder ähnlichen Programmen weitgehend per Hand durchgeführt. Ziel dieser Pipeline ist die vollständige Automatisierung dieses Prozesses auf Basis von Open-Source-Werkzeugen.

Die fehlende Interoperabilität zwischen BIM-Authoring-Tools und Analyse-Software ist ein seit Jahren bekanntes Forschungsproblem: Jin et al.~\cite{jin2019bim} identifizieren in ihrer bibliometrischen Übersichtsarbeit Interoperabilität, Semantik und die Lücke zwischen „As-designed"- und „As-built"-Leistung als zentrale Hemmnisse der BIM-gestützten Gebäudeanalyse. Aktuelle Systemati­sierungsarbeiten zeigen, dass die Kopplung von BIM und KI im AECO-Bereich an Dynamik gewinnt — Alves et al.~\cite{alves2025bim} kartieren 14 BIM- und 16 KI-Kompetenzen als Grundlage für smarte Bauprojekte —, der Übergang vom statischen Modell zu einem semantisch angereicherten, abfragbaren Wissensgraphen aber noch weitgehend offen ist~\cite{peng2024kg}. Wang et al.~\cite{wang2024kg} demonstrieren, dass ein wissens\-graphbasiertes Framework auf Basis geometrischer Relationen und HLK-Topologie in der Lage ist, Gebäudeenergie­modelle vollautomatisch aus BIM-Daten zu erzeugen; eine Folgestudie zeigt, dass derartige Pipelines auch bei unvollständigen BIM-Daten über einen hierarchischen Regel\-satz zur HLK-Topologie-Vervollständigung robust bleiben~\cite{wang2026hvac} — ein Konzept, das die vorliegende Pipeline auf den FEM-Bereich der Tragwerks\-analyse überträgt. Semantische Digital-Twin-Architekturen für Mauerwerksbrücken~\cite{chacon2026semantic} sowie IFC-gestützte Synthese annotierter Punktwolken für das Deep-Learning-basierte Semantiksegmentierungstraining~\cite{zhai2022bim} bestätigen den übergreifenden Trend zur semantischen Datenanreicherung im Bauwesen.

Zusätzlich zur reinen Geometrieverarbeitung wird die \textbf{semantische Struktur des Gebäudes} als gerichteter Graph (Scene Graph) modelliert. Dieser Graph dient als Wissensrepräsentation für automatische Schlussfolgerungen über Bauteilbeziehungen, Randbedingungen und Lastpfade und ist direkt als Kontext für Large Language Models (LLMs) nutzbar~\cite{poux2025scene}. Den Eingang des Systems kann künftig eine vollautomatische Scan-to-BIM-Pipeline bilden: Chen \& Luo~\cite{chen2026scan} identifizieren in ihrer Dekadenschau (2014–2024) die Automatisierung der Geometrierekonstruktion aus Laserscans als zentrales offenes Forschungsfeld — ein direkter Anknüpfungspunkt für Prio~6 der Roadmap.

\subsection{Einordnung in den Projektkontext}

\begin{itemize}[noitemsep]
  \item \textbf{Vorgänger:} \texttt{IFC\_FEM\_3D.ipynb} — BBox-Netz, manuelle Randbedingungen, erste funktionierende IFC-to-JaxFEM-Pipeline
  \item \textbf{Dieses Notebook:} Echte IFC-Geometrie, automatischer Scene Graph, automatische BCs
  \item \textbf{Nächste Stufe:} Heterogenes Materialfeld, LLM-Auto-Call, CesiumJS-Ergebnisrückfluss
\end{itemize}

% ═══════════════════════════════════════════════════════════════════════
\section{Pipeline-Architektur}

\begin{figure}[H]
\centering
\begin{minipage}{0.92\textwidth}
\begin{lstlisting}[language={},style={},
  backgroundcolor=\color{code-bg},
  basicstyle=\ttfamily\footnotesize,
  frame=single, numbers=none,
  rulecolor=\color{hsm-gray}]
IFC-Datei (IfcOpenShell)
   +-- Bauteil-Parsing (IfcWall, IfcColumn, IfcSlab, ...)
   |       +-- BBox + Schwerpunkt (Weltkoordinaten, m)
   |       +-- Material (E, nu aus Pset_MaterialMechanical)
   +-- IFC-Relationship-Analyse
   |       +-- IfcRelAggregates       -> contains / part_of
   |       +-- IfcRelContainedInSpatialStructure -> hosted_in
   |       +-- IfcRelAssociatesMaterial -> shared_material
   |       +-- BBox-Spatial           -> above/adjacent/near
   +-- DBSCAN-Clustering (pro Label)
   +-- Feature-Extraktion (Volumen, Hoehe, Kompaktheit)
   +-- NetworkX DiGraph (Scene Graph)
   |       +-- Graphanalyse -> BC + Lastflaechen
   |       +-- Material-Mapping (volumengew. E_eff)
   +-- GMSH Tet4-Netz (echte IFC-Geometrie)
   +-- JaxFEM Solver (lineare 3D-Elastizitaet)
   +-- Ergebnisse -> Scene Graph -> LLM-Prompt
\end{lstlisting}
\end{minipage}
\caption{Vollständige Pipeline \texttt{IFC\_SceneGraph\_FEM.ipynb}}
\end{figure}

\textbf{Einheiten:} kN und m durchgängig (E-Modul in kN/m², Verschiebungen in mm für Ausgabe).

% ═══════════════════════════════════════════════════════════════════════
\section{Konfigurationsparameter}

Alle Parameter sind in Schritt~1 des Notebooks zentral konfigurierbar:

\begin{table}[H]
\centering
\caption{Konfigurationsparameter}
\begin{tabular}{lrl}
\toprule
\textbf{Parameter} & \textbf{Standardwert} & \textbf{Bedeutung} \\
\midrule
\texttt{USE\_DEMO\_MODEL} & \texttt{False} & Demo-IFC oder eigene Datei \\
\texttt{MESH\_SIZE} & 1{,}0\,m & Tet4-Elementgröße \\
\texttt{E\_DEFAULT} & 30\,000 N/m² & Fallback E-Modul (C25/30) \\
\texttt{NU\_DEFAULT} & 0{,}2 & Fallback Querdehnzahl \\
\texttt{GAMMA\_BETON} & 25{,}0\,kN/m³ & Wichte Stahlbeton \\
\texttt{Q\_NUTZLAST} & 5{,}0\,kN/m² & Charakteristische Nutzlast \\
\texttt{DBSCAN\_EPS} & 0{,}6\,m & Clustering-Radius \\
\texttt{DBSCAN\_MIN\_SAMPLES} & 5 & Mindestpunkte pro Cluster \\
\texttt{REL\_THRESHOLD} & 5{,}0\,m & Max. Distanz für Graphkante \\
\texttt{ADJ\_TOLERANCE} & 0{,}05\,m & BBox-Berührungs-Toleranz \\
\bottomrule
\end{tabular}
\end{table}

% ═══════════════════════════════════════════════════════════════════════
\section{Scene Graph — Formale Beschreibung}

\subsection{Graphstruktur}

Der Scene Graph $G = (V, E)$ ist ein gerichteter Graph (NetworkX \texttt{DiGraph}).

\textbf{Knotenmenge $V$:} Jeder Knoten $v_i \in V$ repräsentiert eine Bauteilinstanz (nach DBSCAN-Clustering) mit den Attributen:
\[
  v_i = \bigl(\,\text{label},\; \text{IFC-Name},\; \mathbf{x}_c,\; V_i,\; h_i,\; E_i,\; \nu_i,\; \ldots\bigr)
\]

\textbf{Kantenmenge $E$:} Jede Kante $e_{ij} \in E$ trägt einen Beziehungstyp $r$ und eine Quelle $s \in \{\text{ifc}, \text{geometric}\}$:

\begin{table}[H]
\centering
\caption{Kanten-Typen im Scene Graph}
\begin{tabular}{lll}
\toprule
\textbf{Typ} & \textbf{Quelle} & \textbf{Bedeutung} \\
\midrule
\texttt{above / below} & geometrisch & Vertikale Lagebeziehung ($\Delta z > 0{,}5$\,m) \\
\texttt{adjacent} & geometrisch & BBox-Abstand $< \varepsilon_\text{adj}$ \\
\texttt{near} & geometrisch & Abstand $< d_\text{rel}$ \\
\texttt{inside / contains} & geometrisch & BBox-Einschachtelung \\
\texttt{part\_of / contains} & IFC & \texttt{IfcRelAggregates} \\
\texttt{hosted\_in} & IFC & \texttt{IfcRelContainedInSpatialStructure} \\
\texttt{shared\_material} & IFC & \texttt{IfcRelAssociatesMaterial} \\
\bottomrule
\end{tabular}
\end{table}

\subsection{Automatische Ableitung von Randbedingungen}

Die Graphanalyse leitet aus der Topologie automatisch FEM-Randbedingungen ab:

\begin{table}[H]
\centering
\caption{Graph-Topologie → FEM-Randbedingungen}
\begin{tabular}{ll}
\toprule
\textbf{Graph-Eigenschaft} & \textbf{FEM-Ableitung} \\
\midrule
\texttt{wall}/\texttt{column} ohne eingehende \texttt{below}-Kante & Einspannung am Fuß ($z = z_\text{min}$) \\
\texttt{slab} mit ausgehender \texttt{above}-Kante & Lastfläche oben \\
\texttt{slab} auf höchster $z$-Ebene & Primäre Lastauflagefläche \\
\bottomrule
\end{tabular}
\end{table}

% ═══════════════════════════════════════════════════════════════════════
\section{FEM-Formulierung}

\subsection{Konstitutives Gesetz — Lineare 3D-Elastizität}

Das Materialverhalten wird durch das isotrope Hooke'sche Gesetz beschrieben:

\begin{equation}
  \boldsymbol{\sigma} = \lambda\,\mathrm{tr}(\boldsymbol{\varepsilon})\,\mathbf{I} + 2\mu\,\boldsymbol{\varepsilon}
\end{equation}

mit dem symmetrischen Verzerrungstensor:
\begin{equation}
  \boldsymbol{\varepsilon} = \frac{1}{2}\left(\nabla\mathbf{u} + (\nabla\mathbf{u})^\top\right)
\end{equation}

und den Lamé-Parametern:
\begin{equation}
  \lambda = \frac{E\,\nu}{(1+\nu)(1-2\nu)}, \qquad
  \mu = \frac{E}{2(1+\nu)}
\end{equation}

\subsection{Effektives Material (Volumenhomogenisierung)}

Das effektive Material des Gesamtsystems wird volumengewichtet berechnet (Voigt-Schranke / iso-Dehnungsannahme):
\begin{equation}
  E_\text{eff} = \frac{\sum_i E_i \cdot V_i}{\sum_i V_i}, \qquad
  \nu_\text{eff} = \frac{\sum_i \nu_i \cdot V_i}{\sum_i V_i}
\end{equation}

Diese Formulierung entspricht der oberen Voigt-Schranke der Homogenisierungstheorie und ist für grobe Netzauflösungen mit überwiegend gleichartigen Materialien eine konservative Näherung. Für heterogene Verbundsysteme (z.\,B. Stahl-Beton-Verbund) ist ein bauteilweises Materialfeld anzustreben (vgl.\ Abschnitt~\ref{sec:limits}).

\subsection{Randwertproblem}

Gesucht ist das Verschiebungsfeld $\mathbf{u}: \Omega \to \mathbb{R}^3$ mit:
\begin{align}
  -\mathrm{div}\,\boldsymbol{\sigma} &= \mathbf{0} \quad \text{in } \Omega \\
  \mathbf{u} &= \mathbf{0} \quad \text{auf } \Gamma_D \text{ (Einspannung)} \\
  \boldsymbol{\sigma}\,\mathbf{n} &= \mathbf{q} \quad \text{auf } \Gamma_N \text{ (Flächenlast)}
\end{align}

mit $\mathbf{q} = (0,\,0,\,-q_\text{ges})^\top$ und $q_\text{ges} = \gamma \cdot h + q_\text{Nutz}$.

\subsection{Von-Mises-Vergleichsspannung}

\begin{equation}
  \sigma_\text{vM} = \sqrt{\frac{3}{2}\,\mathbf{s}:\mathbf{s}},
  \qquad \mathbf{s} = \boldsymbol{\sigma} - \frac{1}{3}\,\mathrm{tr}(\boldsymbol{\sigma})\,\mathbf{I}
\end{equation}

\subsection{Diskretisierung — Tet4-Element}

Lineares Tetraederelement mit 4 Knoten und 12 Freiheitsgraden. Die Formfunktions-Ableitungen im physikalischen Raum ergeben sich als:
\begin{equation}
  \frac{\partial N}{\partial \mathbf{x}} = \frac{\partial N}{\partial \boldsymbol{\xi}} \cdot \mathbf{J}^{-1},
  \qquad
  \mathbf{J} = \begin{bmatrix}
    \mathbf{p}_2 - \mathbf{p}_1 \\
    \mathbf{p}_3 - \mathbf{p}_1 \\
    \mathbf{p}_4 - \mathbf{p}_1
  \end{bmatrix}^\top
\end{equation}

% ═══════════════════════════════════════════════════════════════════════
\section{Ergebnisse — Erster Meilenstein}

\subsection{Modell}

Getestet mit: \texttt{Ifc2x3\_Duplex\_Architecture.ifc} (öffentliche IFC-Beispieldatei).

\begin{table}[H]
\centering
\caption{Modell-Kenngrößen}
\begin{tabular}{lr}
\toprule
\textbf{Kenngröße} & \textbf{Wert} \\
\midrule
IFC-Bauteile (gesamt) & 149 \\
Graph-Knoten (nach Clustering) & 13 \\
Graph-Kanten & 24 (IFC: 0 / Geo: 24) \\
Tet4-Elemente & 1\,958 \\
FEM-Knoten & 639 \\
Freiheitsgrade (DOF) & 1\,917 \\
$E_\text{eff}$ & 30\,000\,000\,kN/m² \\
$\nu_\text{eff}$ & 0{,}20 \\
Gesamtlast $q$ & 21{,}81\,kN/m² \\
\bottomrule
\end{tabular}
\end{table}

\subsection{FEM-Ergebnisse}

\begin{table}[H]
\centering
\caption{FEM-Ergebnisse (Meilenstein 1)}
\begin{tabular}{llr}
\toprule
\textbf{Größe} & \textbf{Symbol} & \textbf{Wert} \\
\midrule
Maximale z-Durchbiegung & $u_{z,\text{max}}$ & $-2{,}064$\,mm \\
Maximale Gesamtverschiebung & $|\mathbf{u}|_\text{max}$ & $2{,}065$\,mm \\
Maximale Von-Mises-Spannung & $\sigma_{\text{vM,max}}$ & $7\,529{,}7$\,kN/m² \\
Minimale Druckspannung & $\sigma_{zz,\text{min}}$ & $-9\,072{,}6$\,kN/m² \\
\bottomrule
\end{tabular}
\end{table}

\textbf{Einordnung:} $\sigma_{\text{vM,max}} = 7\,530$\,kN/m² $\approx 7{,}5$\,MPa liegt deutlich unter der charakteristischen Betondruckfestigkeit von $f_{ck} = 25$\,MPa (C25/30). Die Druckspannung $\sigma_{zz,\text{min}} = -9{,}1$\,MPa entspricht etwa $36\%$ von $f_{ck}$ — plausibel für eine Grobnetz-Rechnung mit $h = 1$\,m.

% ═══════════════════════════════════════════════════════════════════════
\section{LLM-Integration}

Der angereicherte Scene Graph wird als strukturierter Textprompt serialisiert, der direkt an Claude oder ein anderes LLM übergeben werden kann:

\begin{lstlisting}[caption={LLM-Prompt-Schema (graph\_to\_llm\_prompt)}]
=== IFC TRAGWERK - SCENE GRAPH + FEM-ERGEBNISSE ===

BAUTEILE (Knoten):
  [wall_0]  IFC=BasicWall:... (IfcWall)  Typ=wall
            Schwerpunkt=(1.2, 0.1, 1.4) m  Volumen=3.136 m3
            E=30000000 kN/m2  uz_min=-1.2 mm  sigma_vM=4200 kN/m2

BEZIEHUNGEN (Kanten):
  wall_0 -> slab_0  [adjacent | geometric]
  slab_0 -> wall_1  [above    | geometric]

KENNGROESSEN:
  Gesamtlast  q  = 21.81 kN/m2
  u_z_max       = -2.064 mm
  sigma_vM_max  = 7529.7 kN/m2
\end{lstlisting}

Dieser Prompt ermöglicht es einem LLM, strukturmechanische Fragen zu beantworten wie:
\begin{itemize}[noitemsep]
  \item „Welches Bauteil ist am stärksten beansprucht?"
  \item „Ist die Durchbiegung normkonform nach DIN EN 1990 ($l/300$)?"
  \item „Welche Maßnahmen reduzieren die kritische Spannung?"
\end{itemize}

% ═══════════════════════════════════════════════════════════════════════
\section{Vergleich mit dem Vorgänger-Notebook}

\begin{table}[H]
\centering
\caption{IFC\_FEM\_3D.ipynb vs. IFC\_SceneGraph\_FEM.ipynb}
\begin{tabular}{lll}
\toprule
\textbf{Merkmal} & \textbf{IFC\_FEM\_3D} & \textbf{IFC\_SceneGraph\_FEM} \\
\midrule
Netzgeometrie & BBox-Tet4 & Echte IFC-Triangulierung \\
Randbedingungen & Manuell & \textbf{Automatisch aus Graph} \\
Lastflächen & Manuell & \textbf{Automatisch aus Topologie} \\
Material & Manuell & \textbf{Automatisch aus IFC} \\
Struktursemantik & Keine & \textbf{DiGraph, 9 Kanten-Typen} \\
LLM-Export & Keine & \textbf{Strukturierter Prompt} \\
Ergebnisrückfluss & Keine & \textbf{FEM $\to$ Graphknoten} \\
\bottomrule
\end{tabular}
\end{table}

% ═══════════════════════════════════════════════════════════════════════
\section{Bekannte Einschränkungen (Version 1.0)}
\label{sec:limits}

\begin{enumerate}[noitemsep]
  \item \textbf{IFC-Kanten:} Im Duplex-Beispiel liefert \texttt{IfcRelAggregates} keine Kanten (0 IFC-Kanten). Alle 24 Kanten sind geometrisch. Das ist modellspezifisch, kein Algorithmusfehler.
  \item \textbf{Elementgröße:} \texttt{MESH\_SIZE = 1{,}0\,m} ist für Produktionsberechnungen zu grob. Empfehlung: 0{,}2–0{,}5\,m.
  \item \textbf{Homogenisierung:} Volumengewichtetes $E_\text{eff}$ — vereinfachend für Verbundsysteme. Bauteilweises $E$-Feld ist anzustreben.
  \item \textbf{LLM-Auto-Call:} Prompt-Ausgabe implementiert, automatischer API-Call noch nicht.
\end{enumerate}

% ═══════════════════════════════════════════════════════════════════════
\section{Nächste Entwicklungsschritte}

\begin{table}[H]
\centering
\caption{Roadmap}
\begin{tabular}{cll}
\toprule
\textbf{Prio} & \textbf{Schritt} & \textbf{Ziel} \\
\midrule
1 & Elementgröße 0{,}3\,m, Konvergenztest & Ergebnisqualität \\
2 & Bauteilweises Materialfeld & Heterogenes Modell \\
3 & IFC4 mit befüllten Psets & Echte Material-BCs \\
4 & LLM-Auto-Call via Claude API & Automatische Bewertung \\
5 & JSON $\to$ CesiumJS-Einfärbung & Visualisierung in geobim.app \\
6 & Scan-to-BIM-Eingang (Punktwolke $\to$ IFC) & Direkteinspeisung aus Bestandserfassung~\cite{chen2026scan, mahmoud2026scan} \\
7 & Gradient-basierte Formoptimierung & JAX-Differenzierbarkeit nutzen \\
\bottomrule
\end{tabular}
\end{table}

% ═══════════════════════════════════════════════════════════════════════
\section{Abhängigkeiten}

\begin{lstlisting}[language={}, numbers=none]
ifcopenshell >= 0.8
jax, jaxlib
numpy
networkx
scikit-learn       # DBSCAN
scipy              # ConvexHull
gmsh
meshio
pyvista
matplotlib
\end{lstlisting}

Conda-Umgebung: \texttt{jaxfem} (siehe \texttt{environment.yml} im JaxFEM-Root).

% ═══════════════════════════════════════════════════════════════════════
\section*{Referenzen}
\addcontentsline{toc}{section}{Referenzen}

\begin{thebibliography}{9}

\bibitem{poux2025scene}
Poux, F. \& Lehtola, V. (2025).
\textit{3D Scene Graphs for Spatial AI and LLMs}.
Data Science Collective.

\bibitem{chacon2026semantic}
Chacón, R., Ramonell, C., Ventosa San Martino, M. \& Posada, H. (2026).
Semantic digital twins for masonry bridges: Structuring geometrical, material and assessment field data.
\textit{Engineering Structures}, 357, 122503.
\url{https://doi.org/10.1016/j.engstruct.2026.122503}

\bibitem{zhai2022bim}
Zhai, R., Zou, J., He, Y. \& Meng, L. (2022).
BIM-driven data augmentation method for semantic segmentation in superpoint-based deep learning network.
\textit{Automation in Construction}, 140, 104373.
\url{https://doi.org/10.1016/j.autcon.2022.104373}

\bibitem{jin2019bim}
Jin, R., Zhong, B., Ma, L., Hashemi, A. \& Ding, L. (2019).
Integrating BIM with building performance analysis in project life-cycle.
\textit{Automation in Construction}, 106, 102861.
\url{https://doi.org/10.1016/j.autcon.2019.102861}

\bibitem{alves2025bim}
Alves, J.~L., Palha, R.~P. \& Almeida~Filho, A.~T. (2025).
Towards an integrative framework for BIM and artificial intelligence capabilities in smart architecture, engineering, construction, and operations projects.
\textit{Automation in Construction}, 174, 106168.
\url{https://doi.org/10.1016/j.autcon.2025.106168}

\bibitem{peng2024kg}
Peng, Y., Au-Yong, C.~P. \& Myeda, N.~E. (2024).
Knowledge graph of building information modelling (BIM) for facilities management (FM).
\textit{Automation in Construction}, 165, 105492.
\url{https://doi.org/10.1016/j.autcon.2024.105492}

\bibitem{wang2024kg}
Wang, M., Lilis, G.~N., Mavrokapnidis, D., Katsigarakis, K., Korolija, I. \& Rovas, D. (2024).
A knowledge graph-based framework to automate the generation of building energy models using geometric relation checking and HVAC topology establishment.
\textit{Energy and Buildings}, 325, 115035.
\url{https://doi.org/10.1016/j.enbuild.2024.115035}

\bibitem{wang2026hvac}
Wang, M., Lilis, G.~N., Mavrokapnidis, D., Katsigarakis, K., Korolija, I. \& Rovas, D. (2026).
Automating HVAC enrichment in building energy models from imperfect BIM data through topology completeness validation.
\textit{Applied Energy}, 404, 127171.
\url{https://doi.org/10.1016/j.apenergy.2025.127171}

\bibitem{chen2026scan}
Chen, Y. \& Luo, X. (2026).
Indoor scan-to-BIM workflows: Progress, challenges, and future directions (2014--2024).
\textit{Automation in Construction}, 181, 106578.
\url{https://doi.org/10.1016/j.autcon.2025.106578}

\bibitem{mahmoud2026scan}
Mahmoud, M., Li, Y., Adham, M., Mansour, A. \& Chen, W. (2026).
Indoor 3D point cloud reconstruction for scan-to-BIM automation.
\textit{Automation in Construction}, 184, 106853.
\url{https://doi.org/10.1016/j.autcon.2026.106853}

\bibitem{jaxfem}
Xue, T. et al.\ JaxFEM: A differentiable finite element solver in JAX.
\url{https://github.com/tianjuxue/jax-fem}

\bibitem{ifcopenshell}
IfcOpenShell contributors.
IfcOpenShell — Open Source IFC Library.
\url{https://ifcopenshell.org}

\end{thebibliography}

\vfill
\noindent\rule{\linewidth}{0.4pt}
\noindent{\footnotesize\color{hsm-gray}
  Tragwerkslabor · Fachbereich Technik · Fachrichtung Bauingenieurwesen · Hochschule Mainz\\
  University of Applied Sciences · \texttt{philipp.schaefer@hs-mainz.de} · April 2026
}

\end{document}
